\chapter{Derived flags and Hall coefficients}
\label{cyf:chapter}

A flag correspondence can be used with coherent sheaves or
with constructible sheaves. Coherent coefficients retain
nilpotent functions. Constructible coefficients retain the
topology and its supported maps. A degeneration of length-three
quotients separates these two calculations. Its complete flags
acquire a projective line, while their generic point has
scheme-theoretic length six. The multiplicities determine a
cycle trace.

The coherent direct image is free of rank six
over the parameter ring. Its map to functions on the special
fiber has a three-dimensional kernel.

\section{A noncurvilinear point and its complete flags}
\label{cyf:section}

Let \(k\) be a field of characteristic zero. A length-three
quotient of the plane need not admit one algebra generator.
This first occurs at
\[
                       A_0=k[x,y]/(x,y)^2.
\]
Every element of \(A_0\) is \(c+w\), with \(c\in k\) and
\(w^2=0\), so it generates a subalgebra
of dimension at most two.

No change of the chosen polynomial
coordinate makes its multiplication cyclic on \(A_0\).
The pair \(x,y\), together with the unit, nevertheless
generates the algebra. It gives a chart in which both the
quotient and its ideal flags can be calculated through this
point.

\begin{definition}
A local length-three \(k\)-algebra with residue field \(k\)
is called \emph{curvilinear} if it is isomorphic to
\(k[u]/(u^3)\).
\end{definition}
Thus \(A_0\) is noncurvilinear.

Failure of a fixed coordinate to generate is weaker.
For instance \(k[x,y]/(x,y^2)\) is not generated by \(x\),
but is generated by \(x+y\). The failure at \(A_0\)
persists under every such change of coordinate.

\subsection{The simultaneous multiplication chart}

Consider the open of \(\operatorname{Hilb}^3(\mathbb A_k^2)\)
where \(1,x,y\) form the tautological basis. Write its
border products as
\[
 x^2=a+bx+cy,\qquad xy=d+ex+fy,\qquad y^2=g+hx+iy.
\]

\begin{theorem}\label{cyf:thm:chart}
This open is \(\mathbb A_k^6\), with coordinates
\(b,c,e,f,h,i\). Its other coefficients are
\begin{equation}\label{cyf:eq:border}
 a=c(e-i)+f(f-b),\qquad d=ch-ef,\qquad
                      g=e(e-i)+h(f-b).
\end{equation}
The universal multiplication matrices, in the basis
\((1,x,y)\), are
\[
 X=\begin{pmatrix}0&a&d\\1&b&e\\0&c&f\end{pmatrix},
 \qquad
 Y=\begin{pmatrix}0&d&g\\0&e&h\\1&f&i\end{pmatrix}.
\]
These formulas represent the open chart over every
commutative \(k\)-algebra.
\end{theorem}

\begin{proof}
The coefficients of \(x,y\) in
\(x^2y=x(xy)\) and \(xy^2=(xy)y\) give
\eqref{cyf:eq:border}. Their remaining constant
coefficients vanish after this substitution. Equivalently,
direct multiplication gives \(XY=YX\).

Over any base algebra, the matrices therefore define a
\(k[x,y]\)-module structure on the free rank-three module.
The vector \(1\) generates it, since \(X1=x\) and
\(Y1=y\). The resulting surjection from the polynomial
algebra has an ideal as kernel and gives a rank-three
quotient with the prescribed basis. Conversely, any quotient
in the open has unique matrices of this form, and
commutativity forces the displayed eliminations. The two
constructions are inverse and commute with base change.
\end{proof}

Set \(c=t\) and \(b=e=f=h=i=0\). The resulting algebra
over \(k[t]\) is
\begin{equation}\label{cyf:eq:family}
 \mathcal A=k[t,x,y]/(x^2-ty,xy,y^2).
\end{equation}
It is free with basis \(1,x,y\). Its fiber at zero is
\(A_0\), whereas for a nonzero value of \(t\) it is
\(k[x]/(x^3)\), with \(y=x^2/t\).
This last presentation belongs only to the punctured base.

\subsection{The full ordinary flag scheme}

Let \(T=\mathbb A_k^1\). Define \(p:\mathscr F\to T\)
to parameterize ideal flags
\[
 0\subset L\subset M\subset\mathcal A,
 \qquad\operatorname{rank}L=1,
 \qquad\operatorname{rank}M=2,
\]
with all successive quotients locally free. The invariance
conditions under multiplication are closed in the relative
flag variety of \(\mathcal A\), so \(p\) is proper.
It is the complete nested Hilbert correspondence of lengths
\(1,2,3\) pulled back along \eqref{cyf:eq:family}.

The quotient \(\mathcal A/M\) has a unital character
\(x\mapsto\lambda\), \(y\mapsto\mu\). Its equations are
\begin{equation}\label{cyf:eq:character}
                  \lambda^2=t\mu,\qquad
                  \lambda\mu=0,\qquad\mu^2=0.
\end{equation}
Its kernel has basis \(x-\lambda,y-\mu\).

A line in it is locally generated by
\(A(x-\lambda)+B(y-\mu)\), with \([A:B]\in\mathbb P^1\).
Multiplication by \(x\) and \(y\) gives the respective
coefficient columns
\[
                (-A\lambda-B\mu,At),\qquad
                (0,-A\lambda-B\mu).
\]
Taking determinants with \((A,B)\) shows that the line is
an ideal precisely when
\begin{equation}\label{cyf:eq:line}
 tA^2+\lambda AB+\mu B^2=0,
 \qquad \lambda A^2+\mu AB=0.
\end{equation}
In particular the quotient character retains nilpotent
values in a family. It has not been fixed to the origin.

\begin{theorem}\label{cyf:thm:flags}
On \(B\ne0\), put \(v=A/B\). The coordinate ring is
\begin{equation}\label{cyf:eq:B-chart}
 R_B=k[t,v,\lambda]/(F,G),\qquad
 F=tv^3,\quad G=\lambda^2+tv\lambda+t^2v^2,
\end{equation}
with \(\mu=-tv^2-v\lambda\). On \(A\ne0\), put
\(r=B/A\). Its coordinate ring is
\begin{equation}\label{cyf:eq:A-chart}
 R_A=k[r,\mu]/(\mu^2),\qquad
                      t=0,\quad\lambda=-r\mu.
\end{equation}
The overlap is \(r=v^{-1}\), \(\mu=-v\lambda\).
\end{theorem}

\begin{proof}
On the first chart the first line equation eliminates
\(\mu\). The second becomes \(-F=0\), and the first
character equation becomes \(G=0\). The other character
equations reduce to
\[
                 \lambda\mu=-vG+tF,\qquad
                 \mu^2=v^2G+\lambda F.
\]
They impose no further ordinary equations.

On the second
chart the second line equation gives \(\lambda=-r\mu\).
The first then gives \(t=0\). The character equations
reduce to \(\mu^2=0\). On the overlap \(v\) is a unit,
so \(tv^3=0\) gives \(t=0\) and the stated transition.
Conversely these substitutions satisfy every equation
\eqref{cyf:eq:character}--\eqref{cyf:eq:line}.
\end{proof}

After inverting \(t\), set \(q=tv\). Then
\begin{equation}\label{cyf:eq:coinvariant}
 R_B[t^{-1}]
 =k[t,t^{-1},q,\lambda]/(q^3,\lambda^2+q\lambda+q^2).
\end{equation}
It is free over \(k[t,t^{-1}]\) with basis
\(1,q,q^2,\lambda,q\lambda,q^2\lambda\).
Every nonzero fiber is one nonreduced point of length six.
The \(A\)-chart is empty there.

At \(t=0\), the two
charts are \(k[v,\lambda]/\lambda^2\) and
\(k[r,\mu]/\mu^2\), glued by the same transition. The
special fiber is a ribbon with reduced curve \(\mathbb P^1\).
Its points are the lines in the square-zero ideal \((x,y)\).

The total reduced scheme is a horizontal copy of \(T\)
and a vertical \(\mathbb P^1\), meeting at one point.
Their generic scheme multiplicities are six and two,
respectively. Although \(\mathcal A\) is flat over
\(k[t]\), its flag scheme is not: \(v^3\ne0\) in
\(R_B\), as its image modulo \(t\) shows, while
\(tv^3=0\).

\begin{proposition}\label{cyf:prop:root-exchange}
The exchange \(q\leftrightarrow\lambda\) defines an
automorphism of \eqref{cyf:eq:coinvariant}. It has no
extension to a \(k[t]\)-algebra endomorphism of \(R_B\)
with these prescribed root images.
\end{proposition}

\begin{proof}
The three successive \(x\)-eigenvalues on the flag are
\(q,-q-\lambda,\lambda\). Their elementary symmetric
functions vanish. In the displayed quotient,
\(\lambda^3=q^3=0\), so interchanging \(q,\lambda\)
preserves its defining ideal and is an involution.
An extension would satisfy \(t\varphi(v)=\lambda\).
But \(\lambda\) has nonzero image in
\(R_B/tR_B=k[v,\lambda]/\lambda^2\), which excludes this
equation.
\end{proof}

There is a corresponding global obstruction on \(\mathscr F\).
Let \(\kappa_L\) be the eigenvalue of \(x\) on the
tautological line \(L\). It is \(tv\) on the \(B\)-chart
and zero on the \(A\)-chart, so
\(\kappa_L\in t\mathcal O_{\mathscr F}\) locally.
Every automorphism over \(T\) preserves that ideal sheaf.
It therefore cannot pull \(\kappa_L\) back to the quotient
eigenvalue \(\lambda\), which is not in the ideal on the
\(B\)-chart. This excludes the prescribed root exchange
on the whole flag scheme as well.

\subsection{The incidence enhancement and its specialization}

We use the cochain grading, with all polynomial variables
in degree zero. Before elimination on the \(B\)-chart,
the five incidence equations are
\[
\begin{aligned}
 r_1&=\lambda^2-t\mu,&r_2&=\lambda\mu,&r_3&=\mu^2,\\
 r_4&=tv^2+v\lambda+\mu,&r_5&=v^2\lambda+v\mu.
\end{aligned}
\]
Their specified derived zero locus has algebra
\begin{equation}\label{cyf:eq:incidence}
 K_{\mathrm{inc}}=k[t,v,\lambda,\mu]
       \otimes\Lambda(e_1,\ldots,e_5),
 \qquad |e_j|=-1,\quad de_j=r_j.
\end{equation}
The full list of equations is part of this definition.

\begin{theorem}\label{cyf:thm:syzygies}
There is a quasi-isomorphism of commutative dg algebras
\begin{equation}\label{cyf:eq:small}
 K_{\mathrm{inc}}\simeq
 \bigl(k[t,v,\lambda]\otimes\Lambda(e_F,e_G),
                  \ de_F=F,\ de_G=G\bigr)
                  \otimes\Lambda(\chi_1,\chi_2),
\end{equation}
where the two \(\chi\)'s have degree \(-1\) and are
closed. Its cohomology is
\(R_B\otimes\Lambda(\chi_1,\chi_2)\).
\end{theorem}

\begin{proof}
Set \(\nu=r_4\), so
\(\mu=\nu-tv^2-v\lambda\). Change the odd generators by
\begin{align*}
 e_\nu&=e_4,&e_G&=e_1+te_\nu,&e_F&=ve_\nu-e_5,\\
 \chi_1&=e_2-\lambda e_\nu+ve_G-te_F,\\
 \chi_2&=e_3-\bigl(\nu-2v(tv+\lambda)\bigr)e_\nu
                         -v^2e_G-\lambda e_F.
\end{align*}
The even change has coefficient one. Ordered as
\((e_4,e_1,e_5,e_2,e_3)\), the odd change is triangular
with diagonal \((1,1,-1,1,1)\). It is therefore invertible.

Substitution gives
\[
\begin{aligned}
 r_1&=G-t\nu,&r_5&=v\nu-F,\\
 r_2&=\lambda\nu-vG+tF,\\
 r_3&=\bigl(\nu-2v(tv+\lambda)\bigr)\nu+v^2G+\lambda F.
\end{aligned}
\]
Thus the new differentials are \(\nu,G,F,0,0\).

The factor \(k[\nu]\otimes\Lambda(e_\nu)\) retracts to
\(k\), with homotopy
\[
 H(\nu^j)=e_\nu\nu^{j-1}\quad(j\geq1),\qquad
 H(1)=H(e_\nu\nu^j)=0.
\]
With the Koszul tensor sign this gives \eqref{cyf:eq:small}.
Finally \(F\) is a nonzero divisor in the polynomial ring,
and \(G\) is monic in \(\lambda\) over \(k[t,v]/(tv^3)\).
They form a regular sequence, so their Koszul complex
resolves \(R_B\).
\end{proof}

At a nonzero value of \(t\), the dimensions of this
cohomology in degrees \(0,-1,-2\) are \(6,12,6\).
At zero the derived fiber is instead
\begin{equation}\label{cyf:eq:derived-fiber}
 K_{\mathrm{inc}}\otimes^{\mathbb L}_{k[t]}k
 \simeq k[v,\lambda]/(\lambda^2)
               \otimes\Lambda(e_F,\chi_1,\chi_2).
\end{equation}
Indeed the free Koszul model computes derived base change.
Its \(F\)-differential becomes zero, while
\(de_G=\lambda^2\) remains regular. The new \(e_F\)
class records the nonflat flag fiber. It is distinct from
the two incidence syzygies.

Fixing the quotient character to the origin gives another
base change. In \eqref{cyf:eq:incidence}, setting
\(\lambda=\mu=0\) leaves \(de_4=tv^2\) and four
closed generators. Its ordinary ring is \(k[t,v]/(tv^2)\).
Thus prescribing a quotient character changes even the
ordinary flag ring, as well as the derived equations.

Over \(\mathbb C\), the equation count also fixes a
constructible normalization. For a polynomial section
\(s=(s_1,\ldots,s_r)\) on a smooth complex \(m\)-fold
\(X\), put \(W=\sum_j\zeta_js_j\) on
\(X\times\mathbb A^r\), with projection \(\pi\).
Let \(j:Z(s)\hookrightarrow X\). Fiber-linear dimensional
reduction~\cite{DavisonCriticalCoHA}*{Theorem~A.1} gives
\begin{equation}\label{cyf:eq:critical-shift}
 R\pi_!\phi_W\mathbb Q_{X\times\mathbb A^r}[m+r-1]
                 \simeq j_*\mathbb Q_{Z(s)}[m-r].
\end{equation}
Here \(\phi_W\) is the unshifted vanishing-cycle functor,
extended by zero from the zero locus of \(W\), with
\(\phi_0[-1]=\mathrm{id}\). The coefficients are pulled
back rational constants.

For \((F,G)\) on \(\mathbb A^3\) the output shift is
\(+1\). For the five equations on \(\mathbb A^4\) it is
\(-1\). Eliminating \((\nu,e_\nu)\) preserves \(m-r\).

The invertible odd change in the proof of
Theorem~\ref{cyf:thm:syzygies} induces its inverse dual
change of cut coordinates. In those coordinates the
potential is \(\zeta_\nu\nu+\zeta_GG+\zeta_FF\).
The two cut coordinates corresponding to \(\chi_1,\chi_2\)
are unused. Deleting them changes \(m-r\).

These are
different normalized critical coefficients on the same
classical space. Their identification with an oriented
Hall coefficient requires a comparison of the specified
obstruction theories.

\subsection{Proper pushforward and the Gysin parameter}

For constructible coefficients take \(k=\mathbb C\) and work
on complex points. All sheaf morphisms below belong to the
bounded derived category of rational constructible sheaves.
Let \(i:\{0\}\hookrightarrow T=\mathbb A^1\).
The flag family has a section given by
\(\lambda=\mu=0\), \(L=\mathbb C y\).
Write \(\mathscr L\) for its tautological ideal line.
On the reduced vertical fiber it is
\(\mathcal O_{\mathbb P^1}(-1)\subset
\mathcal O_{\mathbb P^1}\otimes(\mathbb C x\oplus\mathbb C y)\).

\begin{theorem}\label{cyf:thm:proper}
The proper coefficient object is
\begin{equation}\label{cyf:eq:proper}
 \mathscr P:=Rp_*\mathbb Q_{\mathscr F}
                  \simeq\mathbb Q_T\oplus i_*\mathbb Q[-2].
\end{equation}
The unit and restriction to the horizontal section split
the first summand. Complex orientation, normalized by
\(\int_{\mathbb P^1}-c_1(\mathscr L)=1\), fixes the second.
\end{theorem}

\begin{proof}
The generic fiber has underlying space a point and the
special fiber has underlying space \(\mathbb P^1\).
Proper base change gives \(R^0p_*\mathbb Q=\mathbb Q_T\),
\(R^2p_*\mathbb Q=i_*\mathbb Q\), and no other cohomology
sheaves. The section retracts the unit. Its complement is
supported at zero and has only \(\mathbb Q\) in degree two,
so it is \(i_*\mathbb Q[-2]\). The stated integral
identifies its generator with the oriented top class.
\end{proof}

Purity supplies the Gysin morphism
\[
 \gamma:i_*\mathbb Q[-2]=i_*i^!\mathbb Q_T
                                      \longrightarrow\mathbb Q_T.
\]
Under adjunction it is the identity of
\(i^!\mathbb Q_T=\mathbb Q[-2]\), and is nonzero.

\begin{theorem}\label{cyf:thm:gysin}
Every degree-zero derived morphism
\(\mathscr P\to\mathbb Q_T\) restricting to the identity
off zero has, in the splitting \eqref{cyf:eq:proper},
the unique form
\begin{equation}\label{cyf:eq:gysin-family}
                         f_a=(\mathrm{id},a\gamma),
                         \qquad a\in\mathbb Q.
\end{equation}
All these maps induce the same ordinary stalk and global
cohomology maps. On compactly supported degree-two
cohomology, their matrix is \((1,a)\).
\end{theorem}

\begin{proof}
Connectedness gives
\(\operatorname{Hom}(\mathbb Q_T,\mathbb Q_T)=\mathbb Q\).
Purity and adjunction give
\[
 \operatorname{Hom}(i_*\mathbb Q[-2],\mathbb Q_T)
 =\operatorname{Hom}(\mathbb Q[-2],\mathbb Q[-2])=\mathbb Q.
\]
The open restriction fixes the first coefficient to one.
The second is the unique multiple of \(\gamma\).

The ordinary pullback of \(\gamma\) at zero is a map
\(\mathbb Q[-2]\to\mathbb Q\), which vanishes in
\(D(\mathbb Q)\). The same holds after ordinary global
sections. Away from zero its source vanishes. Compact
sections instead put both source and target in degree two.
The Gysin map sends the oriented point class to the compact
orientation class of the line, and is the identity in these
bases. This gives \((1,a)\) and distinguishes every \(a\).
\end{proof}

The extraordinary costalk detects the same scalar.
Applying \(i^!\) to the counit makes
\(i^!\gamma:\mathbb Q[-2]\to\mathbb Q[-2]\) the
identity, so \(i^!f_a\) also has matrix \((1,a)\).
The costalk ranks alone do not determine this matrix.

The value \(a=0\) is restriction to the section.
The value \(a=1\) adds the oriented exceptional contribution.
Both are defined maps. Determining which map an independent
Hall correspondence supplies requires its compact pairing.
Ordinary stalk tests do not choose it.

\subsection{Independent products and higher compatibility}

Take \(r\) independent copies of \(p\), over \(T^r\).
For \(S\subseteq\{1,\ldots,r\}\), let
\(i_S:T^{r-|S|}\hookrightarrow T^r\) set the coordinates
indexed by \(S\) to zero. Proper K\"unneth and
\eqref{cyf:eq:proper} give
\[
 \mathscr P^{\boxtimes r}
   \simeq\bigoplus_S i_{S,*}\mathbb Q_{T^{r-|S|}}[-2|S|].
\]
Let \(\gamma_S\) be the product of the Gysin maps in those
coordinates and the identities in the others.

\begin{proposition}\label{cyf:prop:products}
Every degree-zero morphism
\(\mathscr P^{\boxtimes r}\to\mathbb Q_{T^r}\)
equal to the identity on \((\mathbb G_m)^r\) is uniquely
specified by coefficients \(c_S\in\mathbb Q\), with
\(c_\varnothing=1\), through
\[
                              f|_S=c_S\gamma_S.
\]
A family of such maps preserves all independent external
products, with singleton coefficients \(a_1,\ldots,a_r\),
if and only if
\begin{equation}\label{cyf:eq:product-normalization}
                         c_S=\prod_{j\in S}a_j.
\end{equation}
The first two-coordinate defect is
\((c_{\{1,2\}}-a_1a_2)\gamma_{\{1,2\}}\).
It is nonzero precisely when its scalar is nonzero.
\end{proposition}

\begin{proof}
Coordinate purity gives
\(i_S^!\mathbb Q_{T^r}=\mathbb Q_{T^{r-|S|}}[-2|S|]\).
Adjunction therefore identifies the derived Hom from each
summand to the target with
\(R\Gamma(T^{r-|S|},\mathbb Q)=\mathbb Q\) in degree zero.
This proves the classification. Tensoring the singleton
maps multiplies exactly the coefficients indexed by \(S\),
so gives \eqref{cyf:eq:product-normalization}.
Conversely every refinement into singleton factors requires
these coefficients.

All shifts are even. Under further regrouping the Gysin
map is the same product of coordinate orientation classes,
which gives the compatible higher comparisons. Equivalently,
the indicated mapping spaces have no higher homotopy groups,
so these equal classes have unique homotopies when needed.
On \(H_c^{2r}\), every summand contributes one oriented
line and \(\gamma_S\) has coefficient one. It detects
the stated defect.
\end{proof}

This proposition concerns external products of the explicit
flag families. To apply its criterion to an interacting flag
correspondence requires a comparison with that external
product, preserving the supported coefficient map.

\section{The cycle trace and coherent algebra of the complete flag}
\label{cyt:section}

The complete flag family of \S\ref{cyf:section} has one
topological point over a nonzero parameter, but that point has
scheme-theoretic length six. Its exceptional projective line
has generic length two. These lengths select a trace once the
fundamental cycle is specified. They also enter a different
calculation: the direct image of the structure sheaf. We compute
both operations, keeping their coefficient categories distinct.

Set \(R=k[t]\), with \(k\) as above.
Retain the proper flag family \(p:\mathscr F\to T=\operatorname{Spec}R\)
of \eqref{cyf:eq:family}. Its two charts are
\begin{equation}\label{cyt:eq:charts}
 R_B=R[v,\lambda]/(tv^3,G),\qquad
 R_A=k[r,\mu]/(\mu^2),\qquad
 G=\lambda^2+tv\lambda+t^2v^2.
\end{equation}
On the overlap,
\(r=v^{-1}\), \(t=0\), and \(\mu=-v\lambda\).

Write \(H\) for the reduced horizontal section and
\(V\simeq\mathbb P^1_k\) for the reduced vertical component.
The two equations in the first chart form a regular sequence.
Indeed, \(tv^3\) is a nonzero divisor in the polynomial ring,
and \(G\) is monic in \(\lambda\) over its quotient.
The second chart is a hypersurface. Thus \(\mathscr F\) is
pure of dimension one.

The generic lengths calculated in
Theorem~\ref{cyf:thm:flags} give its fundamental cycle
\begin{equation}\label{cyt:eq:cycle}
                         [\mathscr F]=6[H]+2[V].
\end{equation}

\subsection{The multiplicities determine a specified trace}

For the topological calculation take \(k=\mathbb C\) and rational
constructible coefficients. Let \(\omega_{\mathscr F}\) and
\(\omega_T\) denote the dualizing complexes. For either closed
inclusion \(u:C\hookrightarrow\mathscr F\), with \(C=H,V\),
restriction, complex orientation and the proper counit give
\[
 \mathbb Q_{\mathscr F}[2]\longrightarrow
 u_*\mathbb Q_C[2]=u_*\omega_C\longrightarrow\omega_{\mathscr F}.
\]
Weight these two maps by six and two. Their sum is the
fundamental-class morphism \(\operatorname{cl}_{[\mathscr F]}\).

Define the cycle trace by
\begin{equation}\label{cyt:eq:trace-definition}
 \operatorname{Tr}_{[\mathscr F]}
 =\left(
 Rp_*\mathbb Q_{\mathscr F}[2]
 \xrightarrow{Rp_*\operatorname{cl}_{[\mathscr F]}}
 Rp_*\omega_{\mathscr F}\longrightarrow
 \omega_T=\mathbb Q_T[2]\right)[-2].
\end{equation}
This construction uses the orientations of the smooth components.
It requires no homology-manifold property of their union.

Use the splitting
\(\mathscr P=Rp_*\mathbb Q_{\mathscr F}
\simeq\mathbb Q_T\oplus i_*\mathbb Q[-2]\)
of Theorem~\ref{cyf:thm:proper}, where
\(i:\{0\}\hookrightarrow T\). Its exceptional generator is
normalized by \(\int_V-c_1(\mathscr L)=1\).
Let \(\gamma:i_*\mathbb Q[-2]\to\mathbb Q_T\) be the
oriented Gysin counit of Theorem~\ref{cyf:thm:gysin}.

\begin{theorem}\label{cyt:thm:cycle-trace}
The trace \eqref{cyt:eq:trace-definition} and its normalization
to the identity on \(T\setminus\{0\}\) are
\begin{equation}\label{cyt:eq:trace}
 \operatorname{Tr}_{[\mathscr F]}=(6\operatorname{id},2\gamma),
 \qquad
 \widehat{\operatorname{Tr}}_{[\mathscr F]}
       =\tfrac16\operatorname{Tr}_{[\mathscr F]}
       =(\operatorname{id},\tfrac13\gamma).
\end{equation}
\end{theorem}

\begin{proof}
The map from \(H\) is restriction to the section followed by
the identity trace of \(H\to T\). It is therefore six times
the projection onto \(\mathbb Q_T\).

The map from \(V\) factors through \(V\to\{0\}\to T\).
Integration sends its oriented degree-two generator to one,
after which the closed embedding gives \(\gamma\).
Its degree-zero class integrates to zero: the integration
map has target \(\mathbb Q[-2]\), so its restriction from
\(\mathbb Q[0]\) vanishes. The coefficient of the exceptional
summand is consequently two. The generic coefficient is six,
which gives the normalization in \eqref{cyt:eq:trace}.
\end{proof}

Thus compactly supported degree-two cohomology gives the row
\((1,1/3)\). The reduced cycle \([H]+[V]\) instead gives
\((1,1)\), and restriction to the section gives \((1,0)\).
The orientation fixes each of these maps once its cycle is chosen.

\begin{corollary}\label{cyt:cor:products}
For \(r\) independent copies of the flag family, the normalized
product trace has coefficient \(3^{-|S|}\) on the summand
indexed by \(S\subseteq\{1,\ldots,r\}\) in
Proposition~\ref{cyf:prop:products}. In particular its two-axis
coefficient is \(1/9\). These coefficients preserve every
regrouping of the independent factors.
\end{corollary}

\begin{proof}
The component choosing \(V\) in positions \(S\) has cycle
multiplicity \(6^{r-|S|}2^{|S|}\). External products of
orientation maps and proper traces give the product Gysin map
\(\gamma_S\). Division by \(6^r\) gives \(3^{-|S|}\).
Fubini for the component traces, or the coefficient criterion
of Proposition~\ref{cyf:prop:products}, proves compatibility
with regrouping.
\end{proof}

\subsection{The vertical ribbon and the coherent boundary map}

Return to the field \(k\). The horizontal flat closure has
coordinate algebra
\begin{equation}\label{cyt:eq:flat}
 B^{\mathrm{flat}}=R[v,\lambda]/(v^3,G).
\end{equation}
It is free over \(R\), with ordered basis
\((1,v,v^2,\lambda,v\lambda,v^2\lambda)\).
Indeed, first divide by the monic polynomial \(G\) in
\(\lambda\), then by \(v^3\). In particular its defining
ideal is \(t\)-saturated and agrees with \((tv^3,G)\)
after inverting \(t\).

This proves
\((tv^3,G):t^\infty=(v^3,G)\).
The flat closure has no part in the \(A\)-chart.

Let \(V_{\mathrm{rib}}=\mathscr F\times_T\{0\}\).
Its two square-zero coordinates glue by \(\mu=-v\lambda\),
so it has a retraction \(\rho:V_{\mathrm{rib}}\to V\) with
\[
 \rho_*\mathcal O_{V_{\mathrm{rib}}}
        =\mathcal O_V\oplus\mathcal O_V(1),
 \qquad \mathcal O_V(1)^2=0.
\]
Put \(\mathcal J=\ker(\mathcal O_{\mathscr F}
\to\mathcal O_{\mathscr F^{\mathrm{flat}}})\).

\begin{lemma}\label{cyt:lem:ribbon}
The ideal \(\mathcal J\) is annihilated by \(t\), and
\begin{equation}\label{cyt:eq:ribbon}
 \mathcal J\simeq\rho^*\mathcal O_V(-3),\qquad
 \rho_*\mathcal J\simeq\mathcal O_V(-3)\oplus\mathcal O_V(-2).
\end{equation}
For the reduced embedding \(\iota:V\hookrightarrow V_{\mathrm{rib}}\),
ordinary restriction instead gives
\(\iota^*\mathcal J\simeq\mathcal O_V(-3)\).
Consequently \(H^0(\mathcal J)=0\) and
\(H^1(\mathcal J)=k^3\).
\end{lemma}

\begin{proof}
On the \(B\)-chart, \(\mathcal J=(v^3)\).
The ring \(R_B\) is free on \(1,\lambda\) over
\(k[t,v]/(tv^3)\), where the annihilator of \(v^3\) is
\((t)\). Thus multiplication by \(v^3\) identifies
\(R_B/(t)\) with \(\mathcal J\) on this chart.
On the \(A\)-chart, \(\mathcal J\) is the whole structure
sheaf.

Its frames are therefore \(e_B=v^3\) and \(e_A=1\),
with transition \(e_B=v^3e_A\). This is the pullback of
\(\mathcal O_V(-3)\). The projection formula for \(\rho\)
gives the direct sum in \eqref{cyt:eq:ribbon}.
Quotienting by the square-zero ideal gives the stated restriction.
The cohomology follows from that of the two line bundles on
\(\mathbb P^1\).
\end{proof}

The exact sequence defining \(\mathcal J\) has connecting map
\(\partial:B^{\mathrm{flat}}\to H^1(\mathcal J)\).
We fix its signs by taking the \(B\)-lift minus the \(A\)-lift
on the overlap. In the \(e_B\)-frame the class of a flat
section \(f\) is represented by \(f/v^3\), with coefficients
reduced modulo \(t\). The Laurent quotient retains exactly
\[
 \zeta_1=[v^{-2}],\qquad \zeta_2=[v^{-1}],\qquad
\zeta_3=[\lambda v^{-1}].
\]

The \(B\)-chart removes nonnegative powers. The \(A\)-chart
removes powers at most \(-3\) in the even component and
at most \(-2\) in the \(\lambda\)-component. It follows that
\begin{equation}\label{cyt:eq:boundary}
 \partial(v)=\zeta_1,\quad \partial(v^2)=\zeta_2,\quad
 \partial(v^2\lambda)=\zeta_3,\quad
 \partial(1)=\partial(\lambda)=\partial(v\lambda)=0.
\end{equation}
Each \(\zeta_i\) is killed by \(t\).

\begin{theorem}\label{cyt:thm:coherent}
The entire coherent direct image is the degree-zero algebra
\begin{equation}\label{cyt:eq:algebra}
 Rp_*\mathcal O_{\mathscr F}=\mathcal B[0],\qquad
 \mathcal B=R\{1,q,b,\lambda,c,d\},
\end{equation}
where its embedding into \(B^{\mathrm{flat}}\) is
\(q=tv\), \(b=tv^2\), \(c=v\lambda\), and
\(d=tv^2\lambda\). There is an exact sequence
\begin{equation}\label{cyt:eq:index}
 0\longrightarrow\mathcal B\longrightarrow B^{\mathrm{flat}}
 \xrightarrow{\partial}(R/(t))^3\longrightarrow0.
\end{equation}
The inclusion has determinant \(t^3\) in the displayed bases.
Besides the unit and commutativity, its nonzero basis products are
\begin{equation}\label{cyt:eq:products}
 q^2=tb,\quad q\lambda=tc,\quad qc=d,\quad b\lambda=d,
 \qquad \lambda^2=-t(b+c),\quad \lambda c=-d.
\end{equation}
\end{theorem}

\begin{proof}
Equation~\eqref{cyt:eq:boundary} is surjective.
The long exact sequence of coherent cohomology therefore
identifies \(p_*\mathcal O_{\mathscr F}\) with its kernel
and makes \(R^1p_*\mathcal O_{\mathscr F}\) zero.
The two affine charts have affine intersection, so their
\v Cech complex has no terms above degree one.

The kernel consists precisely of flat sections whose
coefficients of \(v,v^2,v^2\lambda\) are divisible by \(t\).
This gives \eqref{cyt:eq:algebra} and the diagonal inclusion
matrix \(\operatorname{diag}(1,t,t,1,1,t)\).
Multiplication in \eqref{cyt:eq:flat} now gives
\eqref{cyt:eq:products}. For example,
\(\lambda c=v\lambda^2=-tv^2\lambda=-d\).
All omitted positive-degree products contain \(v^3\).
\end{proof}

\begin{proposition}\label{cyt:prop:specialization}
Ordinary specialization is a surjection
\[
 \mathcal B/t\mathcal B\longrightarrow
 H^0(V_{\mathrm{rib}},\mathcal O)
     =k\{1,\lambda,c\}
\]
with square-zero kernel \(k\{q,b,d\}\).
Its source has dimension six and its target dimension three.
The source is a local Gorenstein algebra: its socle, the
annihilator of its maximal ideal, is the line \(kd\).
Its maximal-ideal filtration has Hilbert function \((1,4,1)\),
and it is not a complete intersection.
The missing three classes are retained by derived specialization:
\begin{equation}\label{cyt:eq:derived-fiber}
 0\longrightarrow H^1(\mathcal J)\longrightarrow
 H^0R\Gamma(\mathscr F,
       \mathcal O_{\mathscr F}\otimes_R^L k)
 \longrightarrow H^0(V_{\mathrm{rib}},\mathcal O)
 \longrightarrow0,
\end{equation}
where the middle term is \(\mathcal B/t\mathcal B\).
\end{proposition}

\begin{proof}
The split ribbon has \(H^0\) of dimension three and
\(H^1=0\). On it \(t=0\), so the three generators
\(q,b,d\) restrict to zero. Formula~\eqref{cyt:eq:products}
proves the assertion about the kernel. This ideal is not
annihilated by every nonconstant generator: \(qc=b\lambda=d\).

In \(\mathcal B/t\mathcal B\), the maximal ideal
\(\mathfrak m=(q,b,\lambda,c,d)\) satisfies
\(\mathfrak m^2=kd\) and \(\mathfrak m^3=0\).
Multiplication on \(\mathfrak m/\mathfrak m^2\), in the
basis \((q,b,\lambda,c)\) and with values measured in \(d\), is
\[
 \begin{pmatrix}
 0&0&0&1\\
 0&0&1&0\\
 0&1&0&-1\\
 1&0&-1&0
 \end{pmatrix}.
\]
Its determinant is one. Hence no nonzero class in
\(\mathfrak m/\mathfrak m^2\) belongs to the socle, which
is exactly \(kd\).

There are nine independent quadratic
relations among four minimal generators, since
\(\dim\operatorname{Sym}^2(k^4)=10\) and
\(\dim\mathfrak m^2=1\). An Artinian complete intersection
of embedding dimension four has four defining relations,
so this algebra is not one. In particular it differs from
\(B^{\mathrm{flat}}/tB^{\mathrm{flat}}
=k[v,\lambda]/(v^3,\lambda^2)\), whose Hilbert function
is \((1,2,2,1)\).

The resolution \([R\xrightarrow{t}R]\) in degrees \(-1,0\)
computes the derived fiber. Its degree-minus-one cohomology
sheaf is \(\operatorname{Ann}(t)=\mathcal J\), by the
chart computation in Lemma~\ref{cyt:lem:ribbon}. Its
degree-zero cohomology sheaf is \(\mathcal O_{V_{\mathrm{rib}}}\).
The truncation triangle starts with \(\mathcal J[1]\).
Since \(H^0(\mathcal J)=0\) and both sheaves have no
cohomology above degree one, that triangle gives
\eqref{cyt:eq:derived-fiber}.

Tensoring the two-term \v Cech
complex with the same free resolution identifies its total
complex with \(R\Gamma(\mathscr F,\mathcal O)\otimes_R^L k\).
Theorem~\ref{cyt:thm:coherent} and freeness of \(\mathcal B\)
give the asserted middle term. Freeness of the direct image
does not imply flatness of \(\mathcal O_{\mathscr F}\) over \(R\).
\end{proof}

\subsection{A residue pairing over the whole parameter line}

Define the \(R\)-linear functional \(\varepsilon:\mathcal B\to R\)
by \(\varepsilon(d)=1\) and by zero on the other five basis
elements. This specifies an algebraic residue. It makes no
identification with the constructible cycle trace.

\begin{definition}
An \(R\)-linear functional \(f:\mathcal B\to R\) is called
a \emph{Frobenius form} if the pairing
\((x,y)\mapsto f(xy)\) is perfect over \(R\).
\end{definition}

\begin{theorem}\label{cyt:thm:frobenius}
The pairing \((x,y)\mapsto\varepsilon(xy)\) is perfect over
\(R\). The basis dual to \((1,q,b,\lambda,c,d)\) is
\((d,b+c,q+\lambda,b,q,1)\). The resulting coproduct is
\begin{align}
 \Delta(1)={}&1\otimes d+d\otimes1
       +q\otimes(b+c)+(b+c)\otimes q\notag\\
       &+b\otimes\lambda+\lambda\otimes b,
       \label{cyt:eq:casimir}\\
 \Delta(x)={}&(x\otimes1)\Delta(1).
       \label{cyt:eq:coproduct}
\end{align}
It is coassociative, has counit \(\varepsilon\), and is a
\(\mathcal B\)-bimodule map. Its handle element is
\(\operatorname{mult}\Delta(1)=6d\), whose square is zero.

Every \(R\)-linear functional defining a perfect Frobenius
pairing is uniquely
\(\varepsilon_u(x)=\varepsilon(ux)\) for a unit
\(u\in\mathcal B^\times\). Its coproduct and handle are
\[
 \Delta_u(1)=(u^{-1}\otimes1)\Delta(1),\qquad
 \operatorname{mult}\Delta_u(1)=6u^{-1}d.
\]
In particular, \(u=1+d\) gives
\(\varepsilon_u(1)=1\) and the same handle \(6d\).
No such Frobenius coproduct preserves the unit.
\end{theorem}

\begin{proof}
The multiplication table gives the Gram matrix
\[
 \begin{pmatrix}
 0&0&0&0&0&1\\
 0&0&0&0&1&0\\
 0&0&0&1&0&0\\
 0&0&1&0&-1&0\\
 0&1&0&-1&0&0\\
 1&0&0&0&0&0
 \end{pmatrix}.
\]
Its determinant is \(-1\), and multiplication by the listed
dual basis gives the identity pairing matrix.

Under this perfect pairing \(\Delta\) is adjoint to
multiplication. Pairing a proposed coassociativity identity
with \(x\otimes y\otimes z\) gives the associativity
identity \((xy)z=x(yz)\), so proves it. Dual-basis evaluation
gives both counit identities. The equality
\(\varepsilon((ax)y)=\varepsilon(x(ay))\) gives
\((a\otimes1)\Delta(1)=(1\otimes a)\Delta(1)\),
and hence the bimodule identity. Multiplying the six dual
basis pairs gives \(6d\), and \(d^2=0\) by
\eqref{cyt:eq:products}.

The perfect pairing identifies
\(\operatorname{Hom}_R(\mathcal B,R)\) with \(\mathcal B\)
as a \(\mathcal B\)-module. Thus every functional is
uniquely \(\varepsilon(u\,\cdot)\), and its pairing is
perfect exactly when multiplication by \(u\) is invertible,
equivalently when \(u\) is a unit. Replacing each dual basis
element by its product with \(u^{-1}\) gives the displayed
formulas. For \(u=1+d\), the inverse is \(1-d\), and
\(d^2=0\) proves the normalization and handle assertions.
Finally, \(\Delta_u(1)\) induces an isomorphism from the
dual of the rank-six module \(\mathcal B\) to \(\mathcal B\).
The tensor \(1\otimes1\) has image of rank one, so cannot
equal \(\Delta_u(1)\).
\end{proof}

These operations give the usual commutative Frobenius sewing
over \(R\). The closed genus-\(g\) value
\(\varepsilon((6d)^g)\) is \(0\) for \(g=0\), \(6\)
for \(g=1\), and zero for \(g\geq2\).
For the normalized form \(\varepsilon_{1+d}\), these values
are \(1,6,0,0,\ldots\).

Normalization is therefore a choice
of Frobenius form. It does not construct a physical vacuum
expectation. The coproduct obstruction holds for every such
choice. The ordinary trace of multiplication is different:
it is six on \(1\) and zero on \(q,b,\lambda,c,d\).
Indeed these last elements have positive degrees \(1,2,1,2,3\)
in a grading with \(\deg t=0\), so their multiplication
matrices have zero diagonal. That trace pairing is degenerate.

\subsection{The enhancement remains part of the operation}

\begin{proposition}\label{cyt:prop:enhancements}
The intrinsic obstruction theory given by the cotangent
complex of \(\mathscr F\) has virtual class
\([\mathscr F]=6[H]+2[V]\). For \(k=\mathbb C\), the
normalized trace of this cycle is \((\mathrm{id},\gamma/3)\).
On the \(B\)-chart, the two-equation complete intersection
and the five-equation incidence presentation of
Theorem~\ref{cyf:thm:syzygies} have virtual dimensions one
and minus one, respectively. The usual non-equivariant
virtual class of the latter presentation is zero.
For \(k=\mathbb C\), their rational normalized
linear-potential critical coefficients instead have shifts
\(+1\) and \(-1\), as in
\eqref{cyf:eq:critical-shift}.
\end{proposition}

\begin{proof}
For a regular embedding the normal cone is the normal
bundle. Pulling its class back by the zero section gives
the fundamental cycle. This applies on both complete
intersection charts of \(\mathscr F\). Their intrinsic
cotangent obstruction theories agree on the overlap.
The virtual class of that theory is consequently the cycle
of Theorem~\ref{cyt:thm:cycle-trace}, with its scheme
multiplicities.

The first presentation has three smooth coordinates and
two regular equations. The second has four smooth
coordinates and five equations. Eliminating its regular
pair leaves the same two-equation Koszul algebra together
with two closed generators of degree minus one.
The resulting perfect obstruction theory has virtual
rank \(-1\), so its virtual class belongs to the zero
Chow group \(A_{-1}\) of this ordinary chart.

For a rank-\(r\) linear potential on a smooth
\(m\)-dimensional base, fiber-linear dimensional reduction
gives output shift \(m-r\) in the normalization
\eqref{cyf:eq:critical-shift}. The two values are therefore
\(3-2=1\) and \(4-5=-1\). Neither resulting constant
coefficient complex on the nonempty zero locus is zero.
\end{proof}

The dimensionally reduced critical coefficient is a shifted
constant sheaf on the underlying topological zero locus.
Its object alone does not record the two generic nilpotent lengths.
Choosing the reduced fundamental cycle gives the trace
\((\mathrm{id},\gamma)\). Choosing the scheme cycle gives
\((\mathrm{id},\gamma/3)\) after normalization.
The coefficient sheaf alone selects neither morphism.

The cycle map \eqref{cyt:eq:trace-definition} is consequently
a specified comparison target for an oriented Hall operation.
Identifying them requires that operation's actual coefficient
functor, obstruction theory, orientation and supported map.
Such an identification must also preserve the two unit
insertions and substitution of successive flags. These
conditions compare geometric coefficient maps. The residue
pairing on \(\mathcal B\) does not supply those maps.
